\section{Sudo és jogosultság delegálás}

A \texttt{sudo} lehetővé teszi parancsok futtatását más felhasználó (általában root) nevében. A szabályokat a \texttt{/etc/sudoers} fájl vagy \texttt{visudo} szerkesztőn keresztül kezeljük.

\subsection{Sudoers alapok}
\begin{verbatim}
# Example sudoers line
admin ALL=(root) /usr/sbin/shutdown
# Passwordless command
admin ALL=(root) NOPASSWD:/usr/sbin/reboot
\end{verbatim}

Használat:
\begin{lstlisting}[language=bash]
# Run as root
sudo -u root shutdown -r now
# Teszt
sudo -l  # lists the user's privileges
\end{lstlisting}

Tippek:
\begin{itemize}
  \item Mindig használjuk \texttt{visudo}-t, mert szintaxisellenőrzést végez.
  \item Kerüljük a túl széles NOPASSWD bejegyzéseket — finomítsuk a parancslistát.
  \item Használjunk csoport alapú szabályokat (pl. \%admin) a felügyelet egyszerűsítésére.
\end{itemize}
